\section{Long-circle expansion}\label{L:part3}

\begin{proposition}[Long-circle expansion]\label{L:prop:long-circle}
Fix $\phi\in\{-7\pi/8,-\pi/8\}$. Let \(S_{\rm raw}(N,\phi)\) denote the confluent raw sum just considered. With
\[
 E(y)=y^2/(2H_b),\quad f_0=-\phi/(2\pi),\quad
 {\cal R}=d_{M_b}(a)/d_O(a)=(2k/3)^2/\sqrt{2k}
\]
(the last by the row formulas), define
\[
 F_N(\phi)=
 \frac{S_{\rm raw}(N,\phi)}
 {|w|^N{\cal R}^{N^2}N^{7/16-12E(f_0)}}.
\]
There is a locally finite set \(\Gamma_\phi\subset\mathbb Q_{\ge0}\),
containing \(0\), and fixed polynomials \(P_{\phi,\gamma}\) such that,
for every desired accuracy \(J>0\), one can choose \(A_J>J\) with
\[
 F_N(\phi)=
 \sum_{\substack{\gamma\in\Gamma_\phi\\ \gamma\le A_J}}
 N^{-\gamma}P_{\phi,\gamma}(\log N)+O_{\phi,J}(N^{-J})
\]
for every sufficiently large integer \(N\). The coefficient polynomials
are the same for all truncations and all integers \(N\); no subsequence
or parity choice is made. The polynomial at \(\gamma=0\) is a constant,
which may at this stage vanish. For \(\phi=-\pi/8\), every positive rate
is at least \(1\).

\end{proposition}
\begin{proof}
The estimate has two parts. First, the source confines every root species
except \(-D_2\) to a bounded distance from the seam, with tails stronger
than any exponential. The remaining roots change the kinetic momentum
in one direction on each half-circle. Only boundedly many of the gaps
between them can therefore have low energy. We expand the interactions
across long low-energy gaps and integrate their lengths explicitly. The
main summability issue is to do this without losing the bin-square bound
when the other length integrations are extended.

\paragraph{Normalization and phases.}
Take the Cauchy average for \([p^0]\) on \(|p|=N^{-3/2}\), and write
\[
 \ell=6\log N,\qquad L=2\ell+i\xi,
\]
where \(\xi\) ranges over a fixed bounded interval. Use the common bend
from Section~\ref{L:step6}, supported in large fixed widths about the
seam, with endpoints \(\alpha(\pm\ell)=\pm\xi/2\). Its profile in seam
coordinates is independent of \(\ell\); its slope is small enough for
the stability estimate. The source is at zero, and all root layers keep
their prescribed angular offsets.

In the \(M_b\) expression for \(\mathscr P_N\), the field powers of
\(L\) cancel, including the root measures. The determinant quotient
contributes \(N^{7/16}\), a bounded phase, and a convergent correction
series with orders generated by \(2\) and \(3/2\). The source self
constant contributes \({\cal R}^{N^2}\). We divide the kinetic factor
by \(e^{-L E(f_0)}\), so that every real gap carries the nonnegative
excess rate
\[
 e(y)=E(y)-E(f_0),\qquad y\in f_0+\mathbb Z.
\]
Here \(f_0\) is the unique energy minimizer on its momentum translate.
These factors give the normalization in the proposition.

There are no residual size-dependent signs. Indeed the row identity
\(aM_b^2-aM_b+a=0\) gives
\(aT_{M_b}=a(1-M_b)\), whose contractions with \(D_j\) are
\((0,1,0,0,-2,0,0,1)\). Both this row and \(aT_O\) are integral on
roots, and their sum agrees with \(h\) modulo two. Each mixed
source--root pair in the two order-phase sums thus contributes
\(Nh\gamma\bmod2\), independent of its order. Their sum vanishes by
neutrality and \(ah=0\). Finally \(hOSa/H_b=-1\): the linear Poisson
phase cancels the phase of \(w^N\), and the quadratic phase cancels the
one in the denominator of \(\mathscr P_N\).

\paragraph{Source localization and Taylor expansion.}
For a root at real coordinate \(x\), put \(r'=\ell-|x|\), its distance
from the seam. Write its source profile as
\[
 (aM_b^j\gamma)_{j=0}^5
 =A_1(1,1,0,-1,-1,0)+B_1(0,-1,-1,0,1,1).
\]
For the ordered species
\(\delta,-D_2,-D_1-D_2,-\delta\), the coefficients \(A_1\) all
vanish and \(B_1=(1,0,-1,-1)\). Away from the bend the nearest source
factor has modulus
\[
 \left|\prod_{j=0}^5
 (1-e^{(2\pi i(j-d)-|x|)/6})^{aM_b^j\gamma}\right|^N.
\]
On the opposite side, reverse both indices and angles. For each of the
three nonzero profiles, the chosen layer has \(B_1d<0\), and this
modulus is at most \(\exp(-cNe^{-|x|/6})\). To see the sign, compare
the squared factors at indices \(1,2\) with those at \(4,5\), using
\[
 \cos((d-1)\pi/3)+\cos((d-2)\pi/3)
       =\sqrt3\sin(d\pi/3).
\]
Inside the fixed bend widths the logarithmic costs are bounded, since
\(Ne^{-|x|/6}\) is bounded there. Images cost a bounded factor per
root everywhere. Thus every root with nonzero profile, called
\emph{coupled}, contributes at most
\[
 C\exp(-c e^{r'/6}).
\]
The source self images are bounded. The only uncoupled species is
\(-D_2\), and its axis charge is \(h(-D_2)=1\).

For clarity, the Taylor estimate is for the whole product of source
interactions. Let \(K\) be the root count and
\(R_* =\max(0,\max_{\rm coupled}r')\). Root stability, factorial
integration and the Gaussian tail for \(|f|\gg K+1\) bound the total
absolute normalized volume by \(\exp(C\ell)\), even after reserving a
fixed fraction of the source localization. Hence configurations with
\(R_*\ge\sqrt\ell\) contribute an error smaller than every
\(e^{-A\ell}\).

At fixed seam data, replace \(N^{-1}\) by a complex dummy variable.
The source-coupled frequencies are positive elements \(m'\) of
\(\mathbb Z\pm1/6\), without Jordan terms. A cross path of length
\((2l+1)\ell\pm r'\) has order \((2l+1)6m'-1\); a self-image path
has order \(2l\cdot6m'-2\), \(l\ge1\). These are nonnegative even
integers. Keep the order-zero exponentials unexpanded. On a dummy disk
of radius \(c'e^{-R_*/6}\), higher harmonics cost at most
\(Cc'e^{r'/6}\) per coupled root in the exponent. A sufficiently
small fixed \(c'\) therefore preserves half the localization.
Cauchy's estimate gives, for the resulting source product
\(\mathcal A_N\),
\[
 \mathcal A_N=\sum_{j<M}N^{-2j}a_j+\mathcal E_{M,N},
 \qquad
 |\mathcal E_{M,N}|
 \le C_M N^{-2M} C_*^K
        \exp\left(-c_*\sum_{\rm coupled} e^{r'/6}\right)
\]
on \(R_*\le\sqrt\ell\), for large \(N\). The coefficients satisfy
the same bound without \(N^{-2M}\), with a constant depending on their
order. Here \(C_*,c_*>0\) are fixed independently of that order:
the fixed-order loss \(e^{C_M'R_*}\) is absorbed once by the retained
localization, rather than charged separately to every root. The same
bound permits extending each fixed coefficient to all seam positions.
After integration the Taylor remainder costs at most an additional
\(e^{C\ell}\), so increasing \(M\) gives any required power of
\(N^{-1}\). Expand the determinant correction to a sufficiently high
order as well. All estimates are uniform in the Cauchy variable \(\xi\).

\paragraph{Separating the edge and free lengths.}
Fix a coefficient from these source and determinant expansions. On each
arc, order events by distance from the seam. Place an artificial node a
fixed distance beyond the bend. The \emph{edge} ends at the last of
this node and all coupled roots on that arc; let its length be
\(R_j\), \(j=L,R\). The edge data retain exponential moments of
every fixed order in \(R_j\), by the preceding localization, and their
formulas no longer involve the total length \(\ell\).

Beyond the edge only \(-D_2\) roots remain. Their successive free gaps
\(d_i'\ge0\) satisfy
\[
 \sum_{i\in\mathcal I_j}d_i'=\ell-R_j
       \qquad (j=L,R).
\]
The corresponding momenta \(y_i\in f_0+\mathbb Z\) change by
successive \(+1\)'s or successive \(-1\)'s on an arc. Consequently
\[
 \#\{i\in\mathcal I_j:e(y_i)\le B\}\le C_B,
\]
independently of the root count and of the momentum at the seam.
All imaginary kinetic gaps, order phases and source coefficients depend
only on angles, sorting and edge data. The edge's excess kinetic energy
causes no growth. Thus the remaining length integrals have the form of
sums over edge data and root labels of
\[
 \int_{\substack{d_i'\ge0\\
       \sum_{\mathcal I_j}d_i'=\ell-R_j,\ j=L,R}}
       \mathcal G(\boldsymbol d')
       \prod_i e^{-e(y_i)d_i'}\,d\boldsymbol d',
\]
with one length on each arc eliminated by its sum constraint.
The function \(\mathcal G\) contains the remaining root interactions.
We next expose its dependence on the few long low-energy gaps.

\paragraph{Expanding distant root interactions with a stability margin.}
Choose a large fixed \(s_0\). A free gap is \emph{active} if
\(e(y_i)\le B\) and \(d_i'\ge s_0\). For every directed kernel path
crossing an active gap, including periodic images and self-image paths,
expand its frequency series and then exponentiate. Leave all other
paths intact. Each expansion label adds to an active gap a nonnegative
rate \(\sigma_i\in\tfrac16\mathbb Z_{\ge0}\), counting traversal
multiplicity. Jordan derivatives also supply polynomial powers of gap
lengths. These are the only new dependences on active lengths: every
path length is a sum of whole crossed gaps, fixed edge portions and
fixed imaginary offsets.

The absolute expansion is compatible with stability. Reserve half the
exponential frequency decay on the crossed active gaps. The sum of the
remaining absolute log majorants is at most
\[
 \varepsilon_1\left(1+\sum_I K_I^2\right),
\]
with \(\varepsilon_1\) arbitrarily small when \(s_0\) is sufficiently
large. Indeed these paths have real length at least \(s_0\); sum their
exponentially decaying tails by the unit-size bins of the stability
bound, including the polynomial Jordan factors. The fixed angles and
small bend slopes preserve that decay. Deleting these distant paths
from the intact product costs the same majorant. Both costs are
absorbed by a fixed part of the bin-square penalty.

This estimate remains valid when active lengths vary independently
above \(s_0\). Apply stability on the resulting circle, whose two arcs
need not have equal length, keeping the same seam bend and layer
orders. It also bounds absolute polynomial-coefficient norms when
selected active gaps are reset to \(s_0\): the binomial expansion of a
path's powers is bounded by the same polynomial in its total length.
This shortened-circle estimate is what will justify extending the
remaining length integrals after convolution.

\paragraph{Finite convolution and summable coefficient integrals.}
Choose a retained rate \(b_1\) larger than the desired exponential
accuracy in \(\ell\). Choose \(\lambda\) large enough to dominate the
fixed exponential base for species and ordering choices, and then take
\(A_2\) and \(B\ge A_2\) sufficiently large in terms of \(b_1,\lambda\).
An active gap is \emph{cheap} for a particular expansion label when
\[
 e_i=e(y_i)+\sigma_i\le A_2.
\]
The number of cheap gaps and their total polynomial degree are bounded
in terms of these thresholds, independently of \(K\). Every frequency
that introduces a polynomial power also introduces a positive shift.
The possible cheap rates therefore form a fixed finite separated set.

Suppose both arcs have cheap gaps. For a monomial from the expansion,
integrate the cheap lengths on arc \(j\) first:
\[
 \prod_{i\ {\rm cheap}}(d_i')^{k_i}e^{-e_i d_i'},\qquad
 d_i'\ge s_0,\qquad
 \sum_{i\ {\rm cheap}}d_i'
       =\ell-R_j-\sum_{i\ {\rm other}}d_i'.
\]
Translation of the lower endpoints and elementary convolution, or
partial fractions of the Laplace transform, express this integral on
its feasible interval as a finite exponential-polynomial in the right
hand side. Its rates are among the \(e_i\), and its coefficient bounds
depend only on the fixed thresholds. Keep rates at most \(b_1\), and
extend the edge and other-length integrations in their coefficients to
the full range by dropping feasibility. The retained terms are
polynomials in \(\ell\) times
\(e^{-(e_L+e_R)\ell}\), with coefficients independent of \(\ell\).
Discard configurations with no cheap gap on one of the arcs.

Here are the error and summability bounds for these operations. An
inactive high-energy or active noncheap gap can retain the penalty
\[
 \exp[-(3b_1+2\lambda)d_i'].
\]
For high shifted rates, use the reserved half of the path decay; for
high kinetic rates, reserve part of that energy. The choices of
\(A_2,B\) make the stated penalty available. In applying the absolute
coefficient estimate, reset only the cheap lengths to \(s_0\).
Their count and shifts are bounded, so removing their kinetic factors
there costs a threshold-dependent constant. The shortened-circle
estimate preserves stability of the remaining integrand. It applies
also when the edge and other lengths range beyond the original
feasibility boundary, using their original order and the fixed bend
as a function of seam distance. Source Taylor coefficients depend
only on these edge data and still have their localization bound.
Inactive low-energy gaps have total length at most \(C_Bs_0\).

A discarded convolution rate contributes at most an exponential at
rate \(b_1\) in its feasible total length, with polynomial losses.
Extending a retained coefficient beyond feasibility costs growth at
most \(e^{b_1(R_j+\sum_{\rm other}d_i')}\), again with bounded-degree
powers. In that region those lengths already total at least
\(\ell-O(1)\) on one arc. The edge has exponential penalties of every
order, and all other lengths are either bounded in total or carry the
reserved high penalties. Both errors are therefore arbitrarily small
exponentials in \(\ell\) after the thresholds are chosen. If an arc has
no cheap gap, its sum constraint gives the same conclusion directly;
retain the constraints and eliminate one gap per arc. Cheap lengths on
the other arc cost only polynomial volume.

These estimates are summable over root count. We may restrict
\(|f|\le CK+C_B'\); the complement is arbitrarily exponentially small
by the excess kinetic bound on all lengths and the earlier
\(e^{C\ell}\) volume estimate. Sorting cancels identical-species
factorials and costs a fixed base to the power \(K\). Marking low gaps
and other bounded data costs only a threshold-dependent polynomial in
\(1+K\), which can be absorbed into a fixed extra exponential base and
a threshold-dependent constant. No convolution bound introduces an
accuracy-dependent base per ordinary root. Edge-root positions have
simplex volumes up to \(R_j\), including the outermost position when
it defines that edge, and are summable against their exponential
penalties. Beyond the edges, all but boundedly many integrated gaps
yield a factor at most \(\lambda^{-1}\). Choose \(\lambda\) to beat
the fixed bases. This proves absolute summability of both the extended
coefficient integrals and the stated errors.

\paragraph{Rates, logarithms and compatible truncations.}
The preceding estimates are uniform in \(\xi\), so take the Cauchy
average term by term. Each half-arc contributes an order
\[
 6\bigl(E(y)-E(f_0)\bigr)+n,
       \qquad y\in f_0+\mathbb Z,\quad n\in\mathbb Z_{\ge0}.
\]
Together with source and determinant orders these form a locally finite
set of nonnegative rational rates. Polynomial factors in \(\ell\)
become polynomials in \(\log N\). For \(f_0=1/16\), a positive
momentum half-rate is at least \(7/3\); a positive frequency shift has
order at least \(1\), source corrections start at \(2\), and
determinant corrections at \(3/2\). Thus every positive rate is at
least one in this case.

A zero half-rate requires \(y=f_0\) and no frequency shift. Strict
unit changes make that gap unique on its arc. No polynomial path
factor crosses it, since such a factor would also shift its rate.
Its convolution pole is simple. Hence the zero-rate term has no
logarithm on either arc and is a constant.

Finally, construct two truncations to accuracy higher than any fixed
range of rates and subtract them. A first nonzero power/log term in
their difference could not be absorbed by the higher-order remainder.
Their coefficients therefore agree in that range. This gives the
common family claimed in the proposition, for every sufficiently large
integer \(N\).
\end{proof}
